% GENERATED FILE. Edit canonical lessons, then run npm run generate:books.
% canonical-source-hash: fnv1a64:6d40d85520681bd8
% canonical-lessons: IT-C501-pescare, IT-C501-canna-da-pesca, IT-C501-esca, IT-C501-abboccare, IT-C501-borraccia

\chapter{Un giorno di pesca}
\label{ch:it-un-giorno-di-pesca}
\hlchaptermodality{501}

\begin{chapteropening}
\textbf{By the end of this chapter:} \emph{I can talk about going fishing.}

Complete the last lesson of chapter 501: I can talk about going fishing.
\end{chapteropening}

\section[pescare]{pescare — to fish; to draw (a card)}

\label{lesson:IT-C501-pescare}



\begin{quote}
Before the new one: say the Italian for crochet, then the Italian for to darn.
\end{quote}

\subsection*{You'll want to know: pescare}
\textbf{pescare} — \textquotedblleft{}to fish; to draw (a card)\textquotedblright{}.

\begin{quote}
Ogni domenica mio zio va a pescare al lago.
\end{quote}

\begin{quote}
\emph{Every Sunday my uncle goes fishing at the lake.}
\end{quote}

In card games: \emph{pescare una carta} — to draw a card.

\begin{cousinweb}
\emph{pescare} is inherited from Latin \emph{piscārī} \textquotedblleft{}to fish\textquotedblright{} (EtymDB), from \emph{piscis} \textquotedblleft{}a fish\textquotedblright{}.
\end{cousinweb}

\begin{grammarlens}[title={What you've built}]
One more word for telling, explaining and describing at B1.
\end{grammarlens}

\subsection*{Your turn}
\begin{itemize}
\raggedright
  \item \emph{Say it:} \emph{pescare}
  \item \emph{Say it:} \emph{pescare}, once more
  \item \emph{Say it:} a sentence of your own with \emph{pescare}
  \item \emph{Recall:} say the Italian for crochet, then the Italian for to darn, then say \emph{pescare} again
\end{itemize}

\begin{tcolorbox}[breakable,colback=teal!4,colframe=teal!35!black,title={Before you move on}]
What does pescare mean? (\textquotedblleft{}To fish; to draw (a card)\textquotedblright{}.) Say it once more.
\end{tcolorbox}
\section[la canna da pesca]{la canna da pesca — fishing rod}

\label{lesson:IT-C501-canna-da-pesca}



\begin{quote}
Before the new one: say the Italian for to darn, then the Italian for to fish.
\end{quote}

\subsection*{You'll want to know: la canna da pesca}
\textbf{la canna da pesca} — \textquotedblleft{}fishing rod\textquotedblright{}.

\begin{quote}
Ho comprato una canna da pesca nuova.
\end{quote}

\begin{quote}
\emph{I bought a new fishing rod.}
\end{quote}

\emph{La canna} alone is a reed or a cane.

\begin{cousinweb}
\emph{canna da pesca} is built on \emph{canna} \textquotedblleft{}reed, cane\textquotedblright{}, inherited from Latin \emph{canna} (EtymDB), from Greek \emph{kánna} \textquotedblleft{}a reed\textquotedblright{}, and \emph{pesca} \textquotedblleft{}fishing\textquotedblright{}, from \emph{pescare}, which goes back to Latin \emph{piscis} \textquotedblleft{}a fish\textquotedblright{}.
\end{cousinweb}

\begin{grammarlens}[title={What you've built}]
One more word for telling, explaining and describing at B1.
\end{grammarlens}

\subsection*{Your turn}
\begin{itemize}
\raggedright
  \item \emph{Say it:} \emph{la canna da pesca}
  \item \emph{Say it:} \emph{la canna da pesca}, once more
  \item \emph{Say it:} a sentence of your own with \emph{la canna da pesca}
  \item \emph{Recall:} say the Italian for to darn, then the Italian for to fish, then say \emph{la canna da pesca} again
\end{itemize}

\begin{tcolorbox}[breakable,colback=teal!4,colframe=teal!35!black,title={Before you move on}]
What does la canna da pesca mean? (\textquotedblleft{}Fishing rod\textquotedblright{}.) Say it once more.
\end{tcolorbox}
\section[l'esca]{l'esca — bait}

\label{lesson:IT-C501-esca}



\begin{quote}
Before the new one: say the Italian for to fish, then the Italian for fishing rod.
\end{quote}

\subsection*{You'll want to know: l'esca}
\textbf{l'esca} — \textquotedblleft{}bait\textquotedblright{}.

\begin{quote}
Metto l'esca sull'amo.
\end{quote}

\begin{quote}
\emph{I put the bait on the hook.}
\end{quote}

Also a lure in a wider sense: \emph{fare da esca}.

\begin{cousinweb}
\emph{esca} is inherited from Latin \emph{ēsca} \textquotedblleft{}food, bait\textquotedblright{} (EtymDB), from \emph{edere} \textquotedblleft{}to eat\textquotedblright{} (Lewis \& Short).
\end{cousinweb}

\begin{grammarlens}[title={What you've built}]
One more word for telling, explaining and describing at B1.
\end{grammarlens}

\subsection*{Your turn}
\begin{itemize}
\raggedright
  \item \emph{Say it:} \emph{l'esca}
  \item \emph{Say it:} \emph{l'esca}, once more
  \item \emph{Say it:} a sentence of your own with \emph{l'esca}
  \item \emph{Recall:} say the Italian for to fish, then the Italian for fishing rod, then say \emph{l'esca} again
\end{itemize}

\begin{tcolorbox}[breakable,colback=teal!4,colframe=teal!35!black,title={Before you move on}]
What does l'esca mean? (\textquotedblleft{}Bait\textquotedblright{}.) Say it once more.
\end{tcolorbox}
\section[abboccare]{abboccare — to bite (the bait); to fall for it}

\label{lesson:IT-C501-abboccare}



\begin{quote}
Before the new one: say the Italian for fishing rod, then the Italian for bait.
\end{quote}

\subsection*{You'll want to know: abboccare}
\textbf{abboccare} — \textquotedblleft{}to bite (the bait); to fall for it\textquotedblright{}.

\begin{quote}
Dopo un'ora finalmente un pesce ha abboccato.
\end{quote}

\begin{quote}
\emph{After an hour a fish finally bit.}
\end{quote}

Also of people: \emph{ha abboccato} — he fell for it.

\begin{cousinweb}
\emph{abboccare} is built on \emph{bocca} \textquotedblleft{}mouth\textquotedblright{}, inherited from Latin \emph{bucca} \textquotedblleft{}cheek, mouth\textquotedblright{} (Lewis \& Short).
\end{cousinweb}

\begin{grammarlens}[title={What you've built}]
One more word for telling, explaining and describing at B1.
\end{grammarlens}

\subsection*{Your turn}
\begin{itemize}
\raggedright
  \item \emph{Say it:} \emph{abboccare}
  \item \emph{Say it:} \emph{abboccare}, once more
  \item \emph{Say it:} a sentence of your own with \emph{abboccare}
  \item \emph{Recall:} say the Italian for fishing rod, then the Italian for bait, then say \emph{abboccare} again
\end{itemize}

\begin{tcolorbox}[breakable,colback=teal!4,colframe=teal!35!black,title={Before you move on}]
What does abboccare mean? (\textquotedblleft{}To bite (the bait); to fall for it\textquotedblright{}.) Say it once more.
\end{tcolorbox}
\section[la borraccia]{la borraccia — water bottle, flask}

\label{lesson:IT-C501-borraccia}



\begin{quote}
Before the new one: say the Italian for bait, then the Italian for to bite.
\end{quote}

\subsection*{You'll want to know: la borraccia}
\textbf{la borraccia} — \textquotedblleft{}water bottle, flask\textquotedblright{}.

\begin{quote}
Porto sempre la borraccia piena d'acqua.
\end{quote}

\begin{quote}
\emph{I always carry a full water bottle.}
\end{quote}

A bottle you can fill again, for walks and sport.

\begin{cousinweb}
\emph{borraccia} was borrowed from Spanish \emph{borracha} \textquotedblleft{}a leather wine bottle\textquotedblright{} (EtymDB); the deeper origin is not settled.
\end{cousinweb}

\begin{grammarlens}[title={What you've built}]
One more word for telling, explaining and describing at B1.
\end{grammarlens}

\subsection*{Your turn}
\begin{itemize}
\raggedright
  \item \emph{Say it:} \emph{la borraccia}
  \item \emph{Say it:} \emph{la borraccia}, once more
  \item \emph{Say it:} a sentence of your own with \emph{la borraccia}
  \item \emph{Recall:} say the Italian for bait, then the Italian for to bite, then say \emph{la borraccia} again
\end{itemize}

\begin{tcolorbox}[breakable,colback=teal!4,colframe=teal!35!black,title={Before you move on}]
What does la borraccia mean? (\textquotedblleft{}Water bottle, flask\textquotedblright{}.) Say it once more.
\end{tcolorbox}
